\documentclass[]{article}

\usepackage{amsmath}
\usepackage{multirow}
\usepackage{natbib}
\usepackage{graphicx} 


\begin{document}



In this work, we have addressed the long-standing question of the neutrino mass hierarchy by performing a rigorous Bayesian model comparison between the Normal Hierarchy (NH) and the Inverted Hierarchy (IH). The determination of the hierarchy is a crucial open problem in particle physics, with direct consequences for theories of mass generation and for the experimental search for neutrinoless double-beta decay. Our analysis leverages the unprecedented statistical power of the latest cosmological data to provide a definitive answer.

Our methodology combines the tightest available cosmological constraint on the sum of neutrino masses, $\Sigma m_\nu$, from the Dark Energy Spectroscopic Instrument (DESI) Data Release 2 with the latest global fit to neutrino oscillation data from NuFIT 6.0. We computed the Bayesian evidence for each hierarchy, quantifying their relative probability with the Bayes factor, $K = P(D|\text{NH}) / P(D|\text{IH})$. To ensure the robustness of our conclusions against prior assumptions, we employed two distinct and well-motivated prior frameworks: a physically-motivated hierarchical (SJPV) prior and a conservative, objective (HS) reference prior.

The results provide decisive evidence in favor of the Normal Hierarchy. The DESI DR2 data, within the standard $\Lambda$CDM model, constrains the total neutrino mass to be $\Sigma m_\nu < 0.0642$\,eV at 95\% confidence. This upper limit lies significantly below the minimum mass required by the Inverted Hierarchy ($\Sigma m_\nu \ge 0.099$\,eV), placing the entire IH parameter space in severe tension with observations. This tension translates into a Bayes factor of $K > 460$ even under the most conservative HS prior, a value that corresponds to ``decisive'' evidence on the Jeffreys scale. The preference remains strong ($K > 40$) even when extending the analysis to a more flexible $w_0w_a$CDM cosmological model, demonstrating that the conclusion is not an artifact of the baseline model assumptions.

From these results, we have learned that the combination of state-of-the-art cosmological and oscillation data is sufficient to resolve the neutrino mass hierarchy. The Normal Hierarchy is not just preferred, but the Inverted Hierarchy is now robustly disfavored. This finding allows us to place strong constraints on the absolute neutrino mass scale, with the posterior for $\Sigma m_\nu$ being tightly peaked around its minimum allowed value of $\sim 0.059$\,eV. A direct and significant consequence of this result is a revised outlook for neutrinoless double-beta decay experiments. Our posterior for the effective Majorana mass, $m_{\beta\beta}$, is suppressed into the few-meV range, with a 95\% credible interval of $[0.95, 11.55]$\,meV. This is in stark contrast to the much larger signal predicted by the now-disfavored Inverted Hierarchy and suggests that the detection of this rare decay will be significantly more challenging than previously anticipated, likely lying beyond the sensitivity of the next generation of experiments. In conclusion, the precision of modern cosmology has provided a pivotal answer to a fundamental question in particle physics, decisively favoring the Normal Neutrino Mass Hierarchy and reshaping the landscape for future searches for new physics.

\end{document}
                